%%=============================================================================
%% Conclusie
%%=============================================================================

\chapter{Conclusie}%
\label{ch:conclusie}

% TODO: Trek een duidelijke conclusie, in de vorm van een antwoord op de
% onderzoeksvra(a)g(en). Wat was jouw bijdrage aan het onderzoeksdomein en
% hoe biedt dit meerwaarde aan het vakgebied/doelgroep? 
% Reflecteer kritisch over het resultaat. In Engelse teksten wordt deze sectie
% ``Discussion'' genoemd. Had je deze uitkomst verwacht? Zijn er zaken die nog
% niet duidelijk zijn?
% Heeft het onderzoek geleid tot nieuwe vragen die uitnodigen tot verder 
%onderzoek?

Dit onderzoek had als doel een antwoord te formuleren op de vraag hoe het KYC-proces voor Scrada’s Peppol-platform geautomatiseerd kan worden op een manier die voldoet aan de GDPR, fraudebestendig is en kostenefficiënt blijft. 
Door de ontwikkeling en evaluatie van de Proof-of-Concept (PoC) is aangetoond dat een hybride benadering van mobiele technologie en lokale AI hiervoor de meest haalbare oplossing biedt.

\section{Beantwoording van de onderzoeksvraag en deelvragen}
De beantwoording van de deelvragen vormt de conclusie, met tot slot een samenvattend antwoord op de centrale onderzoeksvraag aan de hand van de drie kernpijlers.:

\subsection{Probleemdomein}
\begin{itemize}
    \item \textbf{Huidig proces:} Het KYC-proces bij Scrada verloopt momenteel manueel: cliënten uploaden kopieën van identiteitsbewijzen, waarna een medewerker deze verifieert via de KBO-databank.
    \item \textbf{Bottlenecks:} Het proces is traag, foutgevoelig bij data-invoer en heeft weinig schaalbaarheid tijdens piekperiodes van onboardings.
    \item \textbf{Stappen van CDD:} Automatisering is toegepast op de klantidentificatie en de automatische matching met externe databronnen.
    \item \textbf{Technologieën:} OCR-extractie, VIES-API koppelingen en web scraping (enkel als demonstratie van de PoC) zijn ingezet voor de automatisering.
    \item \textbf{Data en GDPR:} De huidige opslag van onversleutelde kopieën is niet conform aan de GDPR en opslagtermijnen zijn geminimaliseerd door lokale verwerking.
    \item \textbf{Privacy:} Privacy by Design en Data Sovereignty verlagen risico's door biometrische data lokaal op het toestel te houden.
    \item \textbf{Frauderisico's:} Handmatige verificatie is kwetsbaar voor vervalste ID-\-kopieën zonder biometrische controle.
    \item \textbf{Kosten:} De huidige kosten bestaan uit manuele tijdsgebruik door medewerkers.
\end{itemize}

\subsection{Oplossingsdomein}
\begin{itemize}
    \item \textbf{GDPR-conformiteit:} De eKYC-flow blijft conform via Privacy by Design, waarbij enkel gevalideerde resultaten (zoals de naam van de klant, het bedrijf en biometrie-score) worden gedeeld en biometrische ruwe data direct worden verwijderd.
    \item \textbf{Risico's:} De grootste risico's zijn technologische kwetsbaarheden zoals spoofing, die worden aangepakt via liveness-detectie.
    \item \textbf{Fraude-preventie:} Biometrische liveness-detectie en wiskundige RRN-checksum-validatie maken het proces zeer bestand voor identiteitsfraude.
    \item \textbf{Kostenreductie:} Door de inzet van open-source AI (MobileFaceNet) en on-device verwerking is er geen noodzaak voor dure commerciële SaaS-oplossingen.
\end{itemize}

De automatisering van het KYC-proces wordt gerealiseerd door een verschuiving naar een self-service model. In plaats van een manueel proces waarbij medewerkers documenten controleren, voert de klant de verificatie zelfstandig uit via een mobiele applicatie.
De drie kernpijlers van de onderzoeksvraag zijn als volgt vervuld:

\begin{itemize}
    \item \textbf{Kostenefficiëntie:} Door gebruik te maken van open-source AI-modellen (MobileFaceNet) en on-device verwerking (TensorFlow Lite), is Scrada niet afhankelijk van dure externe SaaS-oplossingen zoals itsme. \\ Bovendien bewijst de PoC dat met technieken zoals webscraping (dat uitdrukkelijk enkel voor de demonstratie van de Poc bedoeld is) 
    en gratis API's (VIES) de instapkosten tot een minimum beperkt blijven. De enige bijkomende kosten die deze applicatie zou hebben, zijn het gebruik van de KBO-API, wat op een schaal als Scrada haalbaar is.
    \item \textbf{Fraudeveiligheid:} De PoC bevat meerdere beveiligingslagen. \\ De biometrische gezichtsvergelijking en actieve liveness-detectie voorkomen identiteitsfraude met foto's of video's. Daarnaast zorgt de wiskundige RRN-checksum-validatie en het heuristische match-algoritme met de KBO-databank voor een duidelijke controle op de authenticiteit van de opgehaalde en ingegeven informatie.
    \item \textbf{GDPR-conformiteit:} Conform aan de privacy-by-design-principes vindt de biometrische analyse volledig lokaal op het toestel plaats. \\ Hierdoor verlaat gevoelige ruwe data (pixels) nooit het apparaat van de klant. Enkel de gevalideerde uitkomsten worden naar de backend verzonden, die al gebruikmaakt van encryptie. Hierdoor zijn de risico's bij gegevensverwerking miniem en conform de Europese privacywetgeving.
\end{itemize}

\section{Reflectie}
De evaluatie van de testresultaten toont aan dat de doelstelling van minimaal 50 procent tijdsbesparing op de totale onboarding behaald wordt, en zelfs meer bij de meeste gevallen (indien er geen fouten optreden). 
\\ Bij manuele verificatie is de gemiddelde doorlooptijd 10 minuten, terwijl de geautomatiseerde flow gemiddeld slechts 2-3 minuten in beslag neemt.
Er is echter een technische trade-off vastgesteld: de rekenintensieve liveness-detectie veroorzaakt een lichte daling in de match-score van de biometrie. 
Voor een productieomgeving zal een nauwkeurige meting van de drempelwaarden (thresholds) essentieel zijn om een balans te vinden tussen veiligheid en gebruiksgemak.

Daarnaast is de huidige kosteloze webscraping-methode voor de KBO-\-gegevens niet aangeraden. Voor een stabiel gebruik van de KBO-databank is het gebruik van de officiële KBO-webservice vereist.

\section{Toekomstvisie}
Hoewel de huidige focus van deze Proof-of-Concept ligt op de directe automatisering via AI en API-integraties, 
biedt de huidige omgeving van Scrada interessante mogelijkheden voor de verdere toekomst. Omdat Scrada al blockchain-technologie inzet voor de integriteit van facturen, 
is het een logische denkoefening om te kijken of deze toepassing ook een meerwaarde kan zijn voor identiteitsbeheer.
In een theoretisch scenario zou de ontwikkelde automatisering kunnen gebruikt worden als de eerste stap in een 'One-Time Verification'-model. 
De gecontroleerde identiteit zou dan als een veilige code op de blockchain bewaard kunnen worden, waardoor een gebruiker zich niet telkens opnieuw hoeft te bewijzen. Hierbij blijft de eindgebruiker de eigenaar van zijn data, terwijl de betrouwbaarheid voor het platform gegarandeerd blijft.
\\ Het is echter belangrijk om te benadrukken dat deze blockchain-integratie voor KYC geen deel uitmaakt van de huidige technische realisatie. De focus van dit onderzoek blijft beperkt tot het leveren van een werkende, kostenefficiënte oplossing die vandaag de dag al ingezet kan worden. 
Tot slot kan in de toekomst ook gekeken worden naar de integratie van andere vormen van identiteitsverificatie, zoals het gebruik van NFC-technologie, die ervoor kan zorgen dat het scannen van de identiteitskaart sneller, gemakkelijker en gebruiksvriendelijker wordt.
Deze PoC vormt als conclusie het fundament voor een snellere en veiligere onboarding binnen het Peppol-netwerk.